% =====================================================================
% MATERIAIS E MÉTODOS (Metodologia)
% Descreva o tipo de pesquisa, materiais, procedimentos e instrumentos.
% Este capítulo demonstra: quadro (tabela com borda), tabela booktabs e
% listagem de código.
% =====================================================================

Lorem ipsum dolor sit amet, consectetur adipiscing elit. Este trabalho adota
abordagem qualitativa, quantitativa ou mista e caracteriza-se como pesquisa
exploratória, descritiva ou explicativa~\parencite{exemplo-livro}. [Descreva o
tipo de pesquisa adotado.]

%-------------------------------------------------------------
\section{Materiais}
\label{sec:materiais}
%-------------------------------------------------------------

O \autoref{quad:materiais} lista os materiais utilizados.

% --- EXEMPLO DE QUADRO (tabela com bordas; legenda acima) ------------
\begin{quadro}[!htb]
  \caption{Materiais e instrumentos utilizados.}
  \label{quad:materiais}
  \centering
  % A fonte fica alinhada à margem esquerda DO QUADRO (não do texto).
  \elementocomfonte{%
  \begin{tabular}{|l|l|c|}
    \hline
    \textbf{Item} & \textbf{Descrição} & \textbf{Quantidade} \\
    \hline
    Item A & Lorem ipsum dolor sit amet & 2 \\
    \hline
    Item B & Consectetur adipiscing elit & 5 \\
    \hline
    Item C & Sed do eiusmod tempor        & 1 \\
    \hline
  \end{tabular}}%
  {Fonte: Elaborado pelo autor (\the\year{}).}
\end{quadro}

A \autoref{tab:dados} apresenta os mesmos dados no estilo \textit{booktabs}
(tabela aberta, sem bordas verticais), recomendado para dados numéricos.

% --- EXEMPLO DE TABELA (estilo booktabs) -----------------------------
\begin{table}[!htb]
  \caption{Exemplo de tabela numérica.}
  \label{tab:dados}
  \centering
  % A fonte fica alinhada à margem esquerda DA TABELA (não do texto).
  \elementocomfonte{%
  \begin{tabular}{lrr}
    \toprule
    \textbf{Categoria} & \textbf{Valor} & \textbf{\%} \\
    \midrule
    Primeira  & 120 & 48,0 \\
    Segunda   &  90 & 36,0 \\
    Terceira  &  40 & 16,0 \\
    \midrule
    \textbf{Total} & \textbf{250} & \textbf{100,0} \\
    \bottomrule
  \end{tabular}}%
  {Fonte: Elaborada pelo autor (\the\year{}).}
\end{table}

%-------------------------------------------------------------
\section{Procedimentos}
\label{sec:procedimentos}
%-------------------------------------------------------------

Os procedimentos foram organizados nas seguintes etapas:

\begin{enumerate}[label=\arabic*.]
  \item \textbf{Etapa um:} lorem ipsum dolor sit amet, consectetur adipiscing;
  \item \textbf{Etapa dois:} sed do eiusmod tempor incididunt ut labore;
  \item \textbf{Etapa três:} ut enim ad minim veniam, quis nostrud exercitation.
\end{enumerate}

O \autoref{lst:exemplo} mostra um exemplo de trecho de código.

% --- EXEMPLO DE LISTAGEM DE CÓDIGO -----------------------------------
\begin{lstlisting}[language=Python, caption={Exemplo de código.}, label={lst:exemplo}]
def media(valores):
    # Calcula a media aritmetica de uma lista
    return sum(valores) / len(valores)

print(media([120, 90, 40]))
\end{lstlisting}
% Fonte do código, alinhada à esquerda e em tamanho reduzido (uniforme com a
% legenda), conforme NBR 14724:2024 (5.8).
\noindent{\ABNTEXfontereduzida Fonte: Elaborado pelo autor (\the\year{}).}
